\section{Introduction}
\label{sec:introduction}

AI capability is not the same object as AI output. A model may contain a stock of algorithms, weights, or knowledge and still require machines to turn that stock into inference or research services. This distinction matters for growth theory. If capability raises the productivity of compute, but capability cannot act without compute, capital accumulation may become a bottleneck. It may also become a second amplifier: more capability raises the return to capital, faster accumulation supplies more research compute, and research compute raises capability again.

This paper isolates that mechanism. There is one reproducible stock of capital, $K$, which can be moved without adjustment costs among conventional production, AI inference, and autonomous AI research. Its dated allocations are denoted by $K_Y$, $K_X$, and $K_R$. They satisfy one resource constraint and are not independent stocks. AI services used in final production are $X=BK_X$, and research services are $BK_R$. Thus $B$ is non-rival but cannot produce either service on its own.

The one-stock benchmark is deliberately stark. Separate sectoral stocks would add installation decisions, sector-specific ownership, relative capital-goods prices, and additional costates. Those margins are useful for studying sluggish reallocation, but they are not needed to answer the prior theoretical question: does requiring capital eliminate the explosive feedback generated by AI research? With a single stock, scarcity is represented exactly once, all uses face the same gross rental rate, and the answer does not depend on arbitrary barriers between physically similar machines.

The central parameter is the elasticity of substitution $\sigma_{XL}$ between effective labor and AI services. The analysis produces three distinct conclusions. If $\sigma_{XL}<1$, the two inputs are gross complements. Along any regular nondegenerate long-run equilibrium, AI services remain tied to effective labor, in the spirit of the labor anchor in semi-endogenous growth models \citep{Jones1995}. Capability can keep growing, but the fractions of aggregate capital used for inference and research vanish. If $\sigma_{XL}=1$, the model admits an exact balanced growth path when the combined production--research feedback is subcritical. If $\sigma_{XL}>1$, the economy can enter an AI-dominated region in which output is asymptotically linear in aggregate capital and proportional to a positive power of capability. Capital then fails to impose an asymptotic ceiling: the regular limiting system contains a positive feedback eigenvalue and reaches a finite-time singularity, a possibility related to the economic-singularity question emphasized by \citet{Nordhaus2021}.

The substitute case requires a precise qualification. We prove that a regular bounded-capability tail is inconsistent with the developer's costate equation and that a regular AI-dominated branch with convergent shares is singular. We do not infer from those results that every admissible initial condition generates a singularity, nor that an infinite-horizon equilibrium cannot exist through an irregular, oscillatory, or boundary path. A trajectory defined only up to its blow-up date is a pre-singular equilibrium branch, not an infinite-horizon equilibrium.

The paper also separates decentralization from technology. Households own capital and the AI developer. Competitive final producers rent conventional capital and buy AI services. The developer holds the non-rival capability stock, rents the capital needed for inference and research, and sells inference services under monopoly pricing. Monopoly rents finance research and implement the shadow value of a non-rival asset that otherwise has no direct rental market. This is one possible ownership structure, not a technological necessity; its advantage is that every payment and rent is explicit.

The order of the argument is intentional. Section~\ref{sec:model} presents technology and resources, then solves households, final producers, and the AI developer, in that order. Only afterwards are markets cleared and equilibrium defined. Section~\ref{sec:regimes} derives the three cases for $\sigma_{XL}$. The appendices collect proofs and curvature conditions.
